% Generated draft: requires mathematical and source review.
\documentclass[11pt]{article}
\usepackage[margin=25mm]{geometry}
\usepackage{fontspec}
\setmainfont{Latin Modern Roman}
\usepackage{amsmath,amssymb,mathtools,unicode-math}
\setmathfont{Latin Modern Math}
\usepackage{enumitem,needspace,xurl,hyperref}
\setlength{\parindent}{0pt}
\setlength{\parskip}{5pt}
\setlength{\emergencystretch}{3em}
\newcommand{\N}{\mathbb{N}}
\newcommand{\Z}{\mathbb{Z}}
\newcommand{\Q}{\mathbb{Q}}
\newcommand{\R}{\mathbb{R}}
\newcommand{\C}{\mathbb{C}}
\newcommand{\F}{\mathbb{F}}
\newcommand{\A}{\mathbb{A}}
\newcommand{\PP}{\mathbb{P}}
\newcommand{\E}{\mathbb{E}}
\newcommand{\eps}{\varepsilon}
\newcommand{\dd}{\,\mathrm{d}}
\begin{document}
\section*{UnsolvedMath ID: 800591}
\textbf{Source problem code:} catalogue entry
\subsection{Definitions and mathematical statement}
The supplied excerpt asks whether \(\mathsf{S}_{\mathrm{fin}}(\mathrm{B}_\Omega,\mathrm{B}_\Omega)\) is consistently additive. It does not define the ambient class of spaces, the cover conventions, or ``additive.'' The following gives a provisional formulation under the usual interpretation involving sets of reals and closure under finite unions; these conventions require confirmation against the complete source.

Let \(X\subseteq\R\) have its subspace topology. For this formulation, a cover of \(X\) is a countable family \(\mathcal U\) of subsets of \(X\) such that \(\bigcup\mathcal U=X\) and \(X\notin\mathcal U\). An \(\omega\)-cover is a cover such that
\[
\forall F\subseteq X\;\bigl(F\text{ finite}\Longrightarrow\exists U\in\mathcal U\;(F\subseteq U)\bigr).
\]
A \(\gamma\)-cover is an infinite cover \(\mathcal U\) such that every \(x\in X\) belongs to all but finitely many members of \(\mathcal U\). Write \(\Omega(X)\) and \(\Gamma(X)\) for the families of open \(\omega\)-covers and open \(\gamma\)-covers of \(X\), respectively. Write \(\mathrm{B}_\Omega(X)\) for the family of \(\omega\)-covers whose members are Borel subsets of \(X\).

For cover families \(\mathcal A(X)\) and \(\mathcal B(X)\), the property \(\mathsf{S}_{\mathrm{fin}}(\mathcal A,\mathcal B)\) means that for every sequence \((\mathcal U_n)_{n\in\N}\) from \(\mathcal A(X)\), there are finite subfamilies \(\mathcal F_n\subseteq\mathcal U_n\) such that
\[
\bigcup_{n\in\N}\mathcal F_n\in\mathcal B(X).
\]
Here finite subfamilies may be empty. The property \(\mathsf{S}_1(\mathcal A,\mathcal B)\) means that for every such sequence there are choices \(U_n\in\mathcal U_n\) satisfying \(\{U_n:n\in\N\}\in\mathcal B(X)\).

Let \(P(X)\) mean that \(X\) satisfies \(\mathsf{S}_{\mathrm{fin}}(\mathrm{B}_\Omega,\mathrm{B}_\Omega)\). Under the provisional interpretation above, the question is whether it is consistent with \(\mathsf{ZFC}\) that
\[
\forall X,Y\subseteq\R\qquad
\bigl(P(X)\land P(Y)\bigr)\Longrightarrow P(X\cup Y).
\]
Thus ``additive'' is interpreted here as preservation under the union of two subsets of the same real line, equivalently under nonempty finite unions.

The excerpt also states: ``The problem is also open for \(\mathsf{S}_1(\Gamma,\Omega)\) and \(\mathsf{S}_{\mathrm{fin}}(\Gamma,\Omega)\).'' Its antecedent is absent from the supplied record. If that sentence refers to consistent additivity, the corresponding questions replace \(P\) in the displayed assertion by each of these two properties. This interpretation of the additional sentence remains unresolved; it is not asserted as an established reading of the source.
\subsection{Short English statement}
Under a provisional sets-of-reals interpretation, ask whether it is consistent with \(\mathsf{ZFC}\) that the union of any two sets satisfying \(\mathsf{S}_{\mathrm{fin}}(\mathrm{B}_\Omega,\mathrm{B}_\Omega)\) again satisfies that property. The domain and meaning of additivity need source confirmation. An additional sentence mentions \(\mathsf{S}_1(\Gamma,\Omega)\) and \(\mathsf{S}_{\mathrm{fin}}(\Gamma,\Omega)\), but its missing antecedent prevents a definitive formulation of those additional questions.
\subsection{Sources}
\begin{itemize}
\item Local UnsolvedMath registry: permanent ID \(800591\), source problem code \texttt{catalogue entry}; title: \emph{Selection Principles in Mathematics --- Problem 17 ([AddQuad])}. The registry attributes the problem to Tsaban, with proposed year 2006, and supplies the status label ``open.'' This label is supplied metadata, not an independently verified mathematical or literature assessment.
\item Supplied source reference: Tsaban, \emph{Selection Principles and Special Sets of Reals: Open Problems} (2006), Problem 17, \url{https://arxiv.org/abs/math/0609601}. The record reports access on 2026-10-16 and extraction by \texttt{arxiv\_latex}; these are supplied provenance details, not a claim of access or verification for this draft.
\item Supplied attribution/context marker: \texttt{[AddQuad]}. The record gives no bibliographic expansion of this marker, so its authorship and precise reference remain unresolved.
\end{itemize}
\end{document}
